%% ma: L11-B15-0060
%% lop: 11
%% chuong: 5
%% bai: 15
%% dang: 11.15.8
%% loai: TN4
%% muc: VD
%% dap_an: "C"
%% nguon: "tai-lieu-hoc-tap-mon-toan-lop-11-theo-chuong-trinh-giao-duc-pho-thong-2018.pdf (Lớp Toán Thầy Duy, Phạm Lê Duy, THPT Chu Văn An), Chương 3 Bài 1, Phần 1, Câu 49 (tr.332)"
%% kien_thuc: "Lớp 11 Bài 15 (giới hạn của dãy số)"
%% trang_thai: cho-duyet
%% ly_do: "Dạng mới đề xuất (11.15.x), chờ giáo viên duyệt dạng; tài liệu gốc không có lời giải mẫu, đáp án do kho tự giải và kiểm bằng sympy."
%% ghi_chu: "Đáp án gốc A (1/4) sai; đúng là C (1/2). Đã kiểm số: S_1=1/3, S_2=3/7, S_3=6/13, ... -> 1/2."
\begin{ex}
Cho dãy số $(u_n)$ như sau $u_n=\dfrac n{1+n^2+n^4}$, $n=1,2,\ldots$ Tính giới hạn $\lim(u_1+u_2+\cdots+u_n)$.
\choice{$\dfrac14$}{$1$}{\True $\dfrac12$}{$\dfrac13$}
\loigiai{\sloppy $1+n^2+n^4=(n^2-n+1)(n^2+n+1)$, nên $u_n=\dfrac12\left(\dfrac1{n^2-n+1}-\dfrac1{n^2+n+1}\right)$. Do $(n-1)^2+(n-1)+1=n^2-n+1$ nên tổng triệt tiêu liên tiếp: $u_1+\cdots+u_n=\dfrac12\left(1-\dfrac1{n^2+n+1}\right)\to\dfrac12$.}
\end{ex}
